\chapter{Capstone：可复现研究审计}
\label{ch:research-audit-capstone}

面对一份陌生的回测报告，最有用的第一步不是判断策略好不好，而是恢复它究竟声称了什么、哪些输入在当时可得、哪些数字可以被独立复核。本章把这些问题放进一个小型案例：从输入和时间开始，经过数学、统计、数值和交易限制，最后写出一张有条件的审计意见。

案例故意把“复现成功”和“研究结论成立”分开。程序可以正确重现三行合成数据，却仍然不能证明真实市场表现、广泛泛化或未来盈利；审计的任务正是把这条边界写清楚。

\section{研究对象与证据边界}

设想你收到一份陌生的回测报告。报告声称：研究者测试了 20 个信号，最终模型通过了多重检验；三笔样本外交易带来 $4\%$ 毛收益，扣除成本后仍有 $3.4\%$。代码可以运行，图表也很完整。

现在暂时不评价策略好坏。先问六个更基本的问题：
\begin{enumerate}[leftmargin=*]
  \item 输入是否就是报告声称的输入，且在决策时真实可得？
  \item 数学目标、样本外边界和多重检验族是否写清并被执行？
  \item 核心数值是否能由独立参考账本或第二算法复核？
  \item 交易成本、市场条件和不可成交状态是否进入结论边界？
  \item 陌生人能否从声明环境复现成功与至少一个预期失败？
  \item 数据、代码、模板和文字是否有可用许可，报告是否保留限制？
\end{enumerate}

复现代码只是起点。代码 review 可以发现实现问题，但还要检查“实现是否计算了正确问题”。一个零缺陷程序仍可能精确复现未来数据泄漏、错误样本外、选择后显著性或不存在的成交。本章将这六个问题串成一次完整复核。

\MFQLead{案例材料与研究边界}

为了回答上面的六个问题，案例只需要四类材料：
\begin{itemize}
  \item \textbf{数据.} 三行合成观测，带事件、可得、决策和目标起点；
  \item \textbf{独立基准.} 解析/枚举参考值、时间条件、统计阈值、成本和数值探针；
  \item \textbf{限制说明.} 明确列出结果不能支持的结论；
  \item \textbf{计算程序.} 一个透明实现，能够重复成功情形并展示至少一个失败情形。
\end{itemize}
\begin{lstlisting}[basicstyle=\ttfamily\scriptsize,breaklines=true]
observation_id,event_at,available_at,decision_at,target_start,gross_return
test-001,2024-07-01T07:30:00+08:00,2024-07-01T08:00:00+08:00,2024-07-01T09:25:00+08:00,2024-07-01T09:30:00+08:00,0.02
test-002,2024-07-02T07:30:00+08:00,2024-07-02T08:00:00+08:00,2024-07-02T09:25:00+08:00,2024-07-02T09:30:00+08:00,-0.01
test-003,2024-07-03T07:30:00+08:00,2024-07-03T08:00:00+08:00,2024-07-03T09:25:00+08:00,2024-07-03T09:30:00+08:00,0.03
\end{lstlisting}
这三行是教学 fixture，不是市场数据；它们只让读者能够从文件内容独立重建本章的时间、收益和成本账本。
\begin{definition}[可复现研究包]
一个可复现研究包是一组版本化输入、代码、环境说明、运行方式、独立基准、反例与限制报告，使陌生读者能够从干净位置重建声明结果。可复现只约束“同一材料得到同一证据”，不自动证明因果、预测、可交易或未来稳定。
\end{definition}

\MFQLead{为什么还要检查报告本身}

设文件字节串为 $D$，密码学哈希为 $h(D)$。审计开始前固定
\[
  h(D)=\texttt{5a2561b90d4e692e1100c84d2689199bd9f008f811a45b442524996290d8c2f0}.
\]
哈希相同只能说明文件没有被改写，并不证明数据真实、无偏或有许可。案例同时检查数据和限制说明：如果把“不能建立泛化”改成“证明广泛泛化”，完整性检查应当提醒读者重新查看这项声明。文件完整性与内容判断是两件不同的事。

\section{输入、时间与目标}

在已有工作目录里运行成功并不等于研究对象已经被恢复，因为缓存、未声明文件或环境变量可能悄悄参与结果。复核时先从材料本身回答：
\begin{enumerate}[leftmargin=*]
  \item 输入是否只包含报告声称的数据和限制说明？
  \item 运行所需的环境、随机源和依赖是否写清？
  \item 数据字段、来源和许可是否与研究问题相符？
  \item 运行结果能否由独立参考账本或第二算法复核？
  \item 改变一项关键条件时，是否会得到可解释的失败？
\end{enumerate}
复核的关键不是目录或命令，而是输入、时间、目标和独立基准之间的关系。

\begin{example}
若删掉临时包中的数据文件，而原仓库仍保留同名文件，正确审计必须失败。若它仍成功，说明程序通过硬编码或当前工作目录偷读了包外状态。“能在我机器上跑”不能替代依赖边界。
\end{example}

\MFQLead{数据与时间审计}

CSV 应包含 observation id、事件时间、可得时间、决策时间、目标起点和毛收益贡献。复核先检查字段约定与三行数量，再逐行验证：
\[
  t_e\le t_a\le t_d<t_y,
  \qquad t_d\in\mathcal C_{2024\text{-}07\text{-}v1},
\]
其中所有时间必须显式带 \texttt{+08:00}，$\mathcal C$ 是版本化教学交易日历。训练截止日 6 月 28 日严格早于测试起点 7 月 1 日，三个决策日都属于最终测试。

完整性检查不能替代时间检查。故障测试把第一行可得时间从 8:00 改为 9:31，并重新计算 CSV 哈希。文件完整性因此仍然成立，但因 $t_a>t_d$，时间检查应报告
\begin{center}
\texttt{timeline gate failed: row 1 was unavailable at decision}。
\end{center}
这个负例很重要：若只改数据而不更新哈希，审计只能证明“文件变了”，不能证明它能识别泄漏。

\section{数学、数值与摩擦限制}

审计者应先读取独立参考账本，再运行代码。假设族含 $m=20$ 项，族错误率目标 $\alpha=0.05$，所以 Bonferroni 阈值为
\[
  \frac{\alpha}{m}=\frac{0.05}{20}=0.0025.
\]
选择结果 $p=0.002$ 通过。这里不是声称该信号“为真”，而是说明它满足声明的保守错误率标准；假设定义、依赖与搜索记录仍需审查。

三行毛收益贡献之和与成本账本为
\[
\begin{aligned}
  R^{\mathrm{gross}}&=0.02-0.01+0.03=0.04,\\
  c_{\mathrm{trade}}&=0.0005+0.0007+0.0008=0.002,\\
  C&=3c_{\mathrm{trade}}=0.006,\\
  R^{\mathrm{net}}&=0.04-0.006=0.034.
\end{aligned}
\]
程序使用 \texttt{math.fsum} 归约收益与成本，参照值来自预先写下的有理小数记录。若程序用同一个函数生成期望值，核对就会变成自我对照。

\MFQLead{结论的逻辑强度}

即使上述数字全部通过，也只能说“给定三行合成数据和固定成本假设，审计程序复现净贡献 $0.034$”。不能说：真实市场存在该收益、20 次包含全部研究搜索、$p$ 值在依赖和非平稳下有效、成本模型可外推、策略有容量。审计报告必须把条件句保留下来。

\MFQLead{数值稳定审计}

本章用 $(10^{16},1,-10^{16})$ 作为消去探针。显式的左到右浮点累加得到 $0$，而更稳定的 \texttt{math.fsum} 得到独立 oracle $1$。审计同时要求两件事：探针仍能暴露朴素算法的问题，稳定算法仍得到 $1$。

不要直接用当前 Python 的内置 \texttt{sum} 代表“朴素求和”。某些解释器版本会改进其浮点算法，导致同一测试跨版本改变。教材显式写出三步累加，目的是验证算法概念而不是某个实现细节。这也是环境声明必须进入证据包的原因。

在真实研究中还应根据对象扩展数值检查：协方差矩阵的条件数、优化解对扰动的敏感性、概率与权重的容差尺度、长时间复利的对数实现、并行归约的确定性。一个消去探针不是“数值稳定认证”，只是一个最小例子。

\MFQLead{摩擦、许可与限制审计}

本章成本只有佣金、价差和冲击三项固定教学值。审计程序要求这三个 key 恰好存在；删除一个或添加未解释项都失败。但固定值没有成交量、参与率、波动率和未成交路径，因此不能产生容量结论。A 股的 T+1、涨跌幅、停牌、复权可得性和融券资格仍按第 16 章条件化协议核验，本数据包不声称实现交易所级撮合。

许可复核要区分四类资产：合成数据、审计代码、原创章节与题解、排版类文件分别遵循各自的许可。每项材料都应说明适用范围；数据许可不自动覆盖书稿，代码许可也不自动覆盖模板。仓库中的许可证清单提供具体路径和标识，正文只保留这个边界。

限制报告有六项，并由哈希冻结。报告必须回答“这个项目不能证明什么”：真实资产表现、广泛泛化、复利收益、校准成本、官方日历和未来盈利。删掉限制会使包看起来更强，却让证据更弱。

\section{证据矩阵与审计意见}

最终意见不是简单的 pass/fail，而是一张“声明---证据---失败---限制”矩阵：
\begin{center}
\begin{tabular}{@{}p{0.18\textwidth}p{0.27\textwidth}p{0.21\textwidth}p{0.18\textwidth}@{}}
\toprule
声明 & 正向证据 & 故障证据 & 限制 \\
\midrule
输入未变 & CSV SHA-256 & 修改后哈希失败 & 不证明数据真实 \\
无前视 & 四时间不等式 & 9:31 可得值被拒绝 & 仅三行 fixture \\
错误率受控 & $0.05/20=0.0025$ & $p=0.003$ 应拒绝 & 搜索族须完整 \\
净贡献 $0.034$ & 独立成本账本 & 缺成本 key 失败 & 固定成本、无容量 \\
数值门禁有效 & $0$ 与 $1$ 探针 & 错误稳定和失败 & 仅覆盖消去 \\
许可与限制存在 & SPDX 与报告哈希 & 删除或改写失败 & 逐资产核验 \\
\bottomrule
\end{tabular}
\end{center}

% Generated by tools/build_figures.py; edit figures/figure-specs.json instead.
\MFQFigure{tex/figures/generated/upper-ch17-audit-matrix.pdf}{研究审计不是单一通过状态，而是数据、数学、数值、摩擦与复现证据的联合矩阵。}{upper-ch17-audit-matrix}





审计意见可分为：通过当前声明；附条件通过；拒绝；证据不足。若核心输入缺失或时间泄漏，不能用“建议以后改进”淡化为附条件通过。若只是样本太小，则可以确认机械复现，同时明确拒绝泛化和绩效声明。

\MFQLead{把复核结果写成一张意见卡}

完整复核应重算三行数据、时间条件、成本和数值探针，并分别展示一条成功路径和几条失败路径。读者需要保留的是证据关系：哪一项观察支持哪一句话，哪一个反例迫使我们收窄结论，而不是某条命令或日志字符串。

把本案例的正向路径写成输出摘要，就是：
\begin{lstlisting}[basicstyle=\ttfamily\scriptsize]
package=upper-capstone-v1
data=research_rows.csv rows=3
gross=0.040000,cost=0.006000,net=0.034000
numeric=passed license=CC0-1.0 licenses=4 limitations=6
\end{lstlisting}
若把第一行的可得时间改为 9:31，时间门禁应失败；若只修改限制文字而不更新其绑定哈希，报告门禁应失败。这样读者可以把“通过”与“为什么会失败”放在同一条证据链中。

\noindent\textbf{延伸阅读。}可复现计算研究的通用工作规范参见\cite{sandve2013reproducible}。

\section*{章末练习}
\section*{口述概念}
\begin{enumerate}[leftmargin=*]
  \item 为什么三笔固定资本贡献可以相加，而连续全仓收益应复利？
  \item 一条数值正确的交易记录，还缺什么才能解释为预测证据？
\end{enumerate}
\section*{笔试推导}
\begin{enumerate}[leftmargin=*]
  \item 逐项计算三笔净贡献与总和。
  \item 当成本乘以 $k$，求盈亏平衡值。
\end{enumerate}
\section*{数值编程}
\begin{enumerate}[leftmargin=*]
  \item 将最后一笔成交比例减半，重算线性成本下的总净贡献。
  \item 将候选数从 20 改成 500，比较给定 $p=0.002$ 的 Bonferroni 判断。
\end{enumerate}
\section*{研究判断}
\begin{enumerate}[leftmargin=*]
  \item 已知事件早晨发生，但数据在 9:31 才发布，能否用于 9:25 决策？
  \item 这三笔合成数据能否支持扩大资金？需要哪些新观察？
\end{enumerate}

\subsection*{分级提示}
先确定资本分母和信息截止，再逐笔乘成交比例。多重检验的候选数来自整个选择过程。扩大资金需要重新研究成交与冲击，不能只把净贡献乘以倍数。


\MFQLead{最终提交物与下册桥接}

提交物包括：复现记录、独立核对页、成功与失败摘要、证据矩阵、限制报告和一页审计意见。每条证据标明文件、操作与观察结果；每条限制标明若要解除还需什么新数据或实验。

完成本章后，你拥有的是“能审查量化研究证据”的上册能力，而不是某条方向的完整模型。进入下册时，六条路线会分别把这些问题带入横截面与中性化、统计套利、机器学习、衍生品、组合风险和事件执行。

\noindent\textbf{本章要点。}审计意见必须把声明、正向证据、故障证据和限制放在同一张表中。完整性哈希、可复现运行和漂亮图表都只是证据的一部分；只有当输入时点、数学目标、统计错误率、数值稳定性和交易条件彼此一致时，结论才值得被保留，而且仍应明确它不能证明什么。
